\section{Open questions}

The five items below are the residual uncertainties that this replication
identified but did not close. Each is stated as a testable question with
its evidentiary basis and a concrete next step that fits inside the
free-endpoint policy.

\begin{enumerate}
\item \textbf{Does the paper's RBE$_{\text{SLD}} = 3.7$ (U2OS) survive a
full Lea--Catcheside joint fit?}
Our direct Eq.~(8) reading gives 2.0--2.3, i.e.\ 25--40\% lower than the
paper's 2.8--3.7. Both come from the same Supp Data 3 spreadsheet, so the
difference lies in the fit ansatz. The abstract's headline
``RBE$_{\text{SLD}} > 2.5$'' is only marginally supported by our reading.
\emph{Next:} implement a full two-parameter ($\mu$, RBE$_{\text{SLD}}$)
joint fit across all inter-fraction intervals via
\texttt{scipy.optimize.least\_squares}; extract uncertainties from the
covariance matrix; compare against both our Eq.~(8) direct reading and
the paper's headline. (uicgpu, free, $<$1~h.)

\item \textbf{Does TOPAS-nBio reproduce Fig.~4's 128.5 DSB/Gy at
129~keV/$\mu$m?}
Our analytical surrogate lands at 216 DSB/Gy (RBE$_{\text{DSB}} \sim
6.2$), a factor of 1.7 above the paper. We attribute this to omitted
ionisation-cluster substructure and SSB-pair saturation, but the paper's
TOPAS-nBio value is not independently verified here.
\emph{Next:} TOPAS-nBio v3.9 (Geant4-DNA) on uicgpu, 2.5~$\mu$m water
nucleus, $\alpha$ at 5.4~MeV, 3.2~nm intra-track SSB-pair criterion,
$\sim 10^4$ histories. Estimated 4--6~h on one A100.

\item \textbf{Does the Medras mechanistic simulator agree with our
analytic cluster-foci fits?}
We reproduce N$_{\text{cl}}$ to $\sim 5\%$ against Fig.~5c--d, but our
model is phenomenological, not Medras. A Medras-vs-cluster-foci
comparison would test whether the cluster picture is emergent from
DSB-misrepair or merely a good fit.
\emph{Next:} clone Medras (McMahon 2021), configure with the paper's
cited parameters ($k_{\text{simple}} = 1.4/h$, $k_{\text{complex}} =
0.16/h$, complex fractions 0.43 X-ray / 0.85 $\alpha$), score foci
kinetics at 5 and 10~Gy for both qualities, compare against our fits and
the paper's Fig.~5. 1--2~d, free.

\item \textbf{Is the additive mixed-field LQ (Eq.~2) formally superior to
alternatives (multiplicative, saturating), or only visually similar?}
Both the paper and our replication show visual agreement of Eq.~2 with
the mixed-field points; neither performs an information-criterion
comparison against a multiplicative or saturating model.
\emph{Next:} compute predicted SF for each mixed-field point under Eq.~2,
multiplicative $S_x \cdot S_\alpha$, $\min(S_x, S_\alpha)$, and a
saturating variant; report AIC per model and an F-test on Eq.~2 vs
multiplicative. Pure \texttt{scipy} on uicgpu; $\sim 1$~h.

\item \textbf{Does floating $\beta_\alpha$ within a plausible range
$(0, 0.1)$~Gy$^{-2}$ change RBE$_{\text{D10}}$ and RBE$_{\text{SLD}}$
meaningfully?}
The paper fixes $\beta_\alpha = 0$ without a sensitivity check; we
inherit that. At 3--4 dose points per cell line, the LQ fit is
under-constrained enough that a small non-zero $\beta_\alpha$ could pull
$\alpha_\alpha$ by 10--15\% and propagate into RBE numbers.
\emph{Next:} refit alpha clonogenic data with $\beta_\alpha$ free; report
its fitted value and $1\sigma$; recompute RBE$_{\text{D10}}$ and, via
Eq.~(8), RBE$_{\text{SLD}}$. If consistent with zero, the paper's
simplification is vindicated; if not, downstream RBEs need re-quoted CIs.
Pure \texttt{scipy.curve\_fit}; $\sim 30$~min.
\end{enumerate}
